%% conclusion.tex
%%

%% ==================
\section{Conclusion and Future Work}
\label{sec:Conclusion}
%% ==================
\subsection{Restatement of Research Question and Objectives}
Modern edge-cloud computing environments are characterized by severe non-stationarity, highly dynamic task arrival patterns, and conflicting operational objectives spanning execution delay, financial cost, carbon emissions, and virtual machine utilization. Standard static scheduling heuristics and unconstrained metaheuristics fail to provide consistent trade-offs across fluctuating workloads. To address this challenge, this study investigated the central research question:

\begin{quote}
\textit{Can an adaptive Multi-Armed Bandit (MAB) hyper-heuristic framework dynamically select optimal low-level scheduling algorithms in real-time to optimize multi-objective execution performance under non-stationary HPC workload conditions, without requiring prior workload profile knowledge?}
\end{quote}

To answer this question, our research pursued three core objectives:
\begin{enumerate}
    \item \textbf{Framework Architecture \& Integration:} Formulate a modular hyper-heuristic scheduling engine that integrates diverse low-level heuristics (list schedulers, local search metaheuristics, and hybrid scheduling techniques) into an active arm pool.
    \item \textbf{Dynamic Reward and Policy Modeling:} Design a multi-objective scalarized reward metric encapsulating Makespan, Cost, Carbon Footprint, and VM Utilization, and evaluate four online MAB action-selection policies ($\epsilon$-Greedy, UCB1, Discounted UCB, and Thompson Sampling) to manage non-stationary regret trade-offs.
    \item \textbf{Trace-Driven Empirical Validation:} Conduct comprehensive empirical evaluations on realistic discrete-event simulation platforms driven by real-world HPC workload traces (HPC2N and NASA Ames iPSC/860) to measure system stability, adaptation trajectories, and policy selection dynamics.
\end{enumerate}

\subsection{Summary of Work Done and Key Findings}
We implemented and evaluated the proposed MAB hyper-heuristic framework across 100 simulation rounds per workload trace. The experimental results demonstrate complete success in addressing the research question and achieving all stated objectives:

\begin{itemize}
    \item \textbf{Superior Trajectory Performance:} On the volatile HPC2N trace, MAB controllers achieved a cumulative makespan reduction of up to $71.4\%$ compared to standalone static heuristics such as \texttt{min\_min}, while consistently outperforming standard list schedulers (\texttt{heft} and \texttt{peft}).
    \item \textbf{Prevention of Catastrophic Degradation:} On the highly concurrent NASA Ames trace, the hyper-heuristic engine successfully identified and suppressed unstable single-heuristic modes (e.g., static \texttt{min\_min}, which accumulated over $1.6 \times 10^9$ seconds in makespan), maintaining bounded makespan growth and steady queue processing across all rounds.
    \item \textbf{Balanced Multi-Objective Trade-Offs:} Normalized radar evaluation profiles confirmed that MAB policies achieved broad Pareto dominance over individual heuristics, attaining optimal outer-bound scores ($> 0.95$) across Makespan, Cost, Carbon Footprint, and VM Utilization on uniform traces, and balanced operational trade-offs ($0.65$--$0.72$) under volatile arrival spikes.
    \item \textbf{Adaptive Selection Evolution:} Policy heatmaps and selection ratio trajectories verified that MAB algorithms dynamically shift selection weight based on trace characteristics—favoring hybrid metaheuristics (\texttt{hghhc}) during volatile bursts on HPC2N, and transitioning toward list schedulers (\texttt{peft}, \texttt{sufferage}) during steady-state processing on NASA Ames.
\end{itemize}

\subsection{Implications, Efficacy, and Methodological Limitations}
\subsubsection{Research Implications and Efficacy}
The primary efficacy of this research lies in proving that \textit{meta-learning via upper-level bandit abstraction} provides an efficient middle ground between fast-but-rigid list heuristics and effective-but-expensive metaheuristics. By operating at the hyper-heuristic decision level, the framework eliminates the need for manual heuristic tuning or offline model pre-training. For cloud infrastructure providers, this translates to automated workload adaptivity, improved SLA compliance, reduced carbon intensity, and minimized energy overhead without sacrificing operational throughput.

\subsubsection{Limitations}
Despite its empirical efficacy, the study has three primary methodological limitations:
\begin{enumerate}
    \item \textbf{Static Memory and Exploration Bounds:} Hyperparameters governing policy exploration (e.g., $c = \sqrt{2}$ in UCB1 and $\gamma = 0.95$ in Discounted UCB) were configured globally. On highly steady traces (such as NASA Ames), unweighted exploration policies (e.g., Thompson Sampling) occasionally incurred unnecessary exploration overhead.
    \item \textbf{Simulation Abstraction and Execution Determinism:} While trace-driven CloudSim Plus environments capture discrete-event scheduling dynamics accurately, they do not fully model non-deterministic real-world phenomena such as hypervisor resource contention, transient network jitter, or container cold-start delays.
    \item \textbf{Coarse Decision Granularity:} Reward evaluation occurs at batch boundaries ($B_k$). When scheduling long-running, heterogeneous DAG batches, delayed reward signals can slow down upper-level policy adjustments.
\end{enumerate}

\subsection{Proposals for Meaningful Future Work}
Rather than simply sweeping platform parameters, follow-up research should focus on structural and architectural advancements:

\subsubsection{Contextual and State-Aware Hyper-Heuristics (LinUCB / Contextual Bandits)}
A natural extension is transitioning from context-free MABs to Contextual Bandits (e.g., LinUCB) or Deep Reinforcement Learning (DRL) hyper-heuristics. By parameterizing the upper-level selector with a real-time state vector
\[
\boldsymbol{x}_k = \left[ \text{CoV}_k,\ \text{Depth}_{\text{DAG}},\ \text{Heterogeneity}_{\text{Task}},\ \text{Utilization}_{\text{VM}} \right],
\]
the framework can map workload features directly to low-level heuristics prior to batch execution, accelerating adaptation in rapidly changing environments.

\subsubsection{Hierarchical and Federated Edge-Cloud Scheduling Architectures}
Future projects can extend this framework to multi-tier, geo-distributed edge-cloud topologies. A hierarchical hyper-heuristic architecture could deploy localized, lightweight MAB agents at edge nodes for fast, low-latency task placement, combined with a global federated aggregator that periodically synchronizes heuristic value estimates ($Q$-values or Beta distributions) across regions without transmitting raw task data.

\subsubsection{Physical Testbed Implementation and Commercialization Potential}
To realize the commercial potential of this research, the hyper-heuristic engine can be productized as a custom control-plane plugin for production orchestrators like Kubernetes (e.g., a custom K8s \texttt{kube-scheduler} or Volcano batch scheduling plugin). Evaluating the engine on a physical multi-region cloud testbed (e.g., AWS EKS or OpenShift) will allow researchers to quantify real-world trade-offs between scheduling computation overhead, API latencies, and actual cloud billing reductions under live industry workloads.